% !TeX root = main.tex
\chapter{Методи та інформаційна технологія виявлення аномальних кодів у фото- та відеооб’єктах веб-ресурсів}
\label{chap:methods_technology}

Сформована в розділі~\ref{chap:models_formalization} модель визначає, які представлення мультимедійного об’єкта потрібно зіставити та яким має бути аналітичний результат, однак сама по собі не встановлює способу одержання цих представлень у робочому процесі. Саме на цьому переході виникає методична проблема: структурні ознаки потрібно виділяти з незмінної байтової послідовності, спектрально-зорові --- з контрольовано декодованих похідних даних, а подієво-поведінкові --- з перебігу оброблення. Якщо виконати ці операції без спільної схеми походження, результати стосуватимуться різних станів об’єкта й не утворюватимуть підстави для одного рішення.

У розділі модель послідовно переводиться в дві методичні процедури та технологічний процес. Метод спектрально-зорового та структурного аналізу формує статичний опис, не змішуючи фізично різні джерела ознак. Метод поведінкового узгодження ознак різного походження встановлює, які події справді належать обробленню того самого об’єкта, і лише після цього інтегрує їх зі статичним описом. Інформаційна технологія задає порядок ізоляції, аналізу, прийняття рішення та реєстрації. Така послідовність потрібна не для формального розмежування термінів, а для збереження причинного зв’язку від початкового об’єкта до кожної підстави підсумкового стану.

\section{Загальна організація процесу виявлення аномальних кодів у фото- та відеооб’єктах веб-ресурсів}
\label{sec:process_organization}

Загальний процес має розв’язати суперечність між необхідністю виконати перетворення та вимогою не втратити початковий стан носія. Декодування потрібне для одержання зорових і часових ознак, проте воно може відкинути невідомі ділянки, нормалізувати службові значення або змінити порядок доступних даних. Реконструкція усуває частину вихідної структури навмисно й тому взагалі не може передувати її перевірці. Звідси випливає перша організаційна вимога: початковий мультимедійний об’єкт $X_{mm}$ зберігається незмінним, а всі операції, здатні змінити його, виконуються над простежуваними похідними представленнями.

Структурний аналіз через це виконується над початковою байтовою послідовністю, а зорові, частотні й часові ознаки формуються з контрольованих похідних даних. До зареєстрованого рішення об’єкт залишається в ізоляції та не передається до публікації. Така вимога водночас зберігає доказове походження структурних відхилень і запобігає використанню ще не перевіреного носія прикладними компонентами веб-системи.

Логіку збереження цих двох джерел свідчень показано на рис.~\ref{fig:process_organization}. Статичний опис $H$ формується з незмінного контейнера та його контрольованих представлень, тоді як причинно пов’язані події $X_b$ надходять окремим каналом. Їх не можна об’єднати раніше, оскільки до перевірки належності журнал може містити фонові події, а статичний опис --- спостереження, одержані за іншої версії засобу. Лише після перевірки походження обидва входи утворюють інтегроване представлення $Z$, з якого формується аналітичний стан $S_A$. П’ять операційних етапів відповідної технології деталізовано на рис.~\ref{fig:information_technology}.

\begin{figure}[H]
\centering
\begin{tikzpicture}[x=1cm,y=1cm,scale=0.93,transform shape]
\node[dissertation small,text width=2.35cm,minimum height=12mm] (object) at (-6.35,1.25)
  {Вхід 1\\Об’єкт $X_{mm}$};
\node[dissertation small,text width=2.75cm,minimum height=12mm] (staticop) at (-3.25,1.25)
  {Формування $X_c$\\
   Статичний аналіз $F_H$};
\node[dissertation small,text width=1.85cm,minimum height=12mm] (profile) at (-0.55,1.25)
  {Результат методу 1\\Опис $H$};

\node[dissertation small,text width=2.35cm,minimum height=12mm] (log) at (-6.35,-1.25)
  {Вхід 2\\Журнал подій $J_b$};
\node[dissertation small,text width=2.75cm,minimum height=12mm] (events) at (-3.25,-1.25)
  {Причинний відбір\\
   $J_b\!\rightarrow\!X_b$ і кодування $F_b$};
\node[dissertation small,text width=1.85cm,minimum height=12mm] (behavior) at (-0.55,-1.25)
  {Проміжний результат\\Ознаки $Z_b$};

\node[dissertation small,text width=2.25cm,minimum height=18mm] (integration) at (2.05,0)
  {Метод 2\\Контроль походження\\
   та доступності\\Інтеграція $F_I$};
\node[dissertation small,text width=2.21cm,minimum height=12mm] (joint) at (4.85,0)
  {Інтегроване\\представлення $Z$};
\node[dissertation small,text width=2.30cm,minimum height=12mm] (state) at (7.65,0)
  {Модельне рішення $F_D$\\
   Вихід: стан $S_A$};

\node[dissertation container,fit=(object)(staticop)(profile),
  label={[font=\footnotesize]above:Метод 1: формування статичного опису}] {};
\node[dissertation container,fit=(log)(events)(behavior),
  label={[font=\footnotesize]below:Метод 2: причинно-фазовий відбір}] {};

\draw[dissertation arrow] (object.east) -- (staticop.west);
\draw[dissertation arrow] (staticop.east) -- (profile.west);
\draw[dissertation arrow] (log.east) -- (events.west);
\draw[dissertation arrow] (events.east) -- (behavior.west);
\draw[dissertation arrow] (profile.east) -- (integration.155);
\draw[dissertation arrow] (behavior.east) -- (integration.205);
\draw[dissertation arrow] (integration.east) -- (joint.west);
\draw[dissertation arrow] (joint.east) -- (state.west);
\end{tikzpicture}
\caption{Роздільне формування статичного представлення та представлення ознак подій і поведінки з подальшим узгодженням}
\label{fig:process_organization}
\end{figure}


Щоб кожний результат можна було пов’язати не лише з об’єктом, а й з конкретним етапом його опрацювання, процес подано як послідовність типізованих станів $S_0,\ldots,S_5$. Кожний наступний стан утворюється з попереднього стану за параметрами відповідного етапу. Таке подання необхідне через те, що зміна правила розбору, параметра перетворення або порога створює інший результат навіть для незмінного початкового об’єкта. Отже, разом зі станом зберігаються версії параметрів, а повторне виконання не підміняє раніше одержаний результат.

Для кожного переходу визначено передумову, очікуваний вихід і стан завершення. Передумова перевіряє належність даних до одного об’єкта та сумісність версій; очікуваний вихід задає мінімальний склад відомостей, потрібний наступному етапу. Стан завершення розрізняє повний результат, явно неповний результат і технічну відмову. Це розмежування принципове, оскільки відсутність спостереження не свідчить про відсутність аномалії. Неповний результат передається далі лише разом із маскою доступності та причиною обмеження, а відмова не перетворюється на нульову ознаку.

З урахуванням цих умов дослідницький перехід від прийнятого носія до відтворюваного рішення організовано в такі етапи:

\begin{enumerate}
\item \textit{Попереднє оброблення.} Об’єкту присвоюється ідентифікатор, обчислюється контрольна сума, зберігається незмінний $X_{mm}$, установлюється фактичний формат, формуються $X_s$ і $X_m$, а контрольоване декодування утворює $X_v$.
\item \textit{Виділення ознак.} Перший метод будує $X_f$ і $X_t$, формує $Z_s$, $Z_v$, локалізації $L_s$, $L_v$ і статичний опис $H$; контрольоване середовище окремо формує $X_b$.
\item \textit{Інтеграція ознак.} Другий метод формує $Z_b$ і виконує перехід $Z=F_I(H,Z_b,m_A,C)$, у якому статичні групи надходять у складі $H$, а маска й контекст явно обмежують допустиме узгодження.
\item \textit{Прийняття рішення.} Функція $F_D$ формує аналітичний стан $S_A$, після чого політика визначає технологічну дію без зміни модельного висновку.
\item \textit{Реєстрація результатів.} Вихідний стан, версії, локалізації, часові межі й посилання на вихідні спостереження зберігаються як один запис.
\end{enumerate}

Однакова глибина перевірки для всіх об’єктів або створює зайві витрати, або змушує скорочувати аналіз саме для складних випадків. Підставою для поглиблення є не лише високе значення однієї статичної оцінки, а й суперечність між гілками: структурно нетиповий, але зорово звичайний об’єкт потребує іншого розбору, ніж об’єкт із двома низькими узгодженими оцінками. Розбіжність гілок характеризується величиною $\delta_A=|\rho_s-\rho_v|$, де $\rho_s$ і $\rho_v$ --- відповідно структурна та спектрально-зорова оцінки, означені в підрозділі~\ref{sec:svs_method}.

Показник $\delta_A$ не визначає загрози --- він указує, що дві гілки по-різному описують стан носія, через що передчасне усереднення приховало б невизначеність. Спочатку перевіряється повнота обов’язкових даних: якщо структурна гілка недоступна або чинна політика потребує відсутнього зорового представлення, об’єкт передається на експертний розгляд. За повних даних поглиблений маршрут призначається, коли принаймні одна оцінка досягає порога $\tau_r$ або розбіжність досягає порога $\tau_d$; в інших випадках достатній базовий маршрут. Базовий маршрут обмежується детермінованими статичними ознаками, тоді як поглиблений додатково зберігає докладніше частотне представлення позначених ділянок і контекст декодування та передбачає контрольоване спостереження оброблення з формуванням $X_b$. Значення $\delta_A$ і обраний маршрут зберігаються в контексті $C$ та враховуються під час обчислення упевненості $c_A$. Пороги $\tau_r$ і $\tau_d$ визначають лише на даних для налаштування та зберігають у версії конфігурації; у контрольних експериментах розділу~\ref{chap:experiments} усі об’єкти опрацьовувалися за повним статичним маршрутом, тому ці пороги характеризують організацію процесу, а не окремо перевірену конфігурацію.

Інформаційний зв’язок етапів узагальнено в таблиці~\ref{tab:ch3_process_stages}. Таблиця фіксує типи входів і результатів кожного етапу та не повторює внутрішніх процедур методів.

\begin{table}[htbp]
\centering
\caption{Входи, основні операції та результати етапів загального процесу виявлення аномальних кодів}
\label{tab:ch3_process_stages}
\small
\begin{tabularx}{\textwidth}{>{\raggedright\arraybackslash}p{0.19\textwidth}YY}
\toprule
\textbf{Етап} & \textbf{Вхід і основна операція} & \textbf{Результат} \\
\midrule
Попереднє оброблення & Прийнятий об’єкт і контекст; ідентифікація, структурний розбір, контрольоване декодування & $X_c$ і початкова маска \\
Виділення ознак & $X_c$ та контрольоване середовище; побудова $X_f$, $X_t$, статичних ознак і причинно пов’язаних подій & $H$, $X_b$ та повне $X_A$ \\
Інтеграція ознак & $Z_v$, $Z_s$, $Z_b$, маска і контекст & $Z$ \\
Прийняття рішення & $Z$, пороги й версійована політика & $S_A$ та технологічна дія \\
Реєстрація результатів & Аналітичний стан, дія, версії та походження & Контрольний запис \\
\bottomrule
\end{tabularx}
\end{table}

Загальна організація, таким чином, відповідає на питання, за яких умов результат одного етапу може стати входом наступного. Вона не утворює третьої моделі або окремого методу, а задає інваріанти, без яких ознаки різного походження не можна обґрунтовано зіставляти. Далі ці інваріанти конкретизовано в операціях двох методів.

\section{Метод спектрально-зорового та структурного аналізу фото- і відеооб’єктів}
\label{sec:svs_method}

Модель розділу~\ref{chap:models_formalization} визначила склад статичного опису, але залишила відкритим питання, як одержати його з контейнера, не втративши ознак, що існують лише до декодування, і не приписавши структурним відхиленням властивостей зорового вмісту. Ця проблема не розв’язується простим об’єднанням усіх вимірювань в один вектор. Додаткова байтова область може не змінювати зображення, а слабка зміна коефіцієнтів може не порушувати граматики формату; отже, нульовий результат однієї гілки не спростовує спостереження іншої.

Метою методу є формування статичного опису ознак можливого прихованого впливу через узгоджене, але роздільне опрацювання двох представлень того самого носія. Структурна гілка аналізує початковий контейнер і службові дані та зберігає байтові межі відхилень. Спектрально-зорова гілка аналізує контрольовано декодоване представлення, його частотні, хвилькові, залишкові й часові характеристики та зберігає просторово-часові локалізації. Об’єднання відбувається лише після формування окремих ознак, оцінок і локалізацій. Через це метод не формує остаточного класу й не визначає технологічної дії: його результат має бути придатним для подальшого причинного зіставлення, а не підміняти його.

Вхід методу становлять передоброблені статичні дані $X_c=\langle X_{mm},\allowbreak X_v,\allowbreak X_s,\allowbreak X_m\rangle$, початкова маска $m_c=(m_s,m_v)$ та контекст $C$; події $X_b$ навмисно вилучено з цього входу, щоб не допустити їх неявного впливу на статичний опис. Представлення $X_f$ і $X_t$ утворюються з $X_v$ уже всередині методу, тому разом із ними зберігаються параметри перетворення та відбору. Виходом є $H=\langle Z_v,\allowbreak Z_s,\allowbreak L_v,\allowbreak L_s,\allowbreak m_c,\allowbreak C\rangle$. Відповідність між цим входом і структурованим результатом установлено формулою~\eqref{eq:ch3_svs_mapping}:
\begin{equation}
H=F_H(X_c,m_c,C,\Theta_H),
\label{eq:ch3_svs_mapping}
\end{equation}
де, $F_H$ --- функція методу спектрально-зорового та структурного аналізу; $X_c$ --- сукупність статичних представлень; $m_c$ --- початкова маска; $C$ --- контекст; $\Theta_H$ --- параметри форматних правил, областей, перетворень, відбору кадрів і оцінювання; $H$ --- статичний опис.

Відображення у формулі~\eqref{eq:ch3_svs_mapping} задає межу відповідальності методу: зміна будь-якого вхідного представлення або параметра створює інший опис і має бути простежуваною. Передумовами виконання є наявність незмінного $X_{mm}$, контрольної суми та підтверджене походження похідних представлень. Після виконання послідовності має бути одержано не просто числові оцінки, а опис із локалізаціями, маскою й контекстом. Цю послідовність визначено так:

\begin{enumerate}
\item перевірити ідентифікатор, контрольну суму, версії засобів і походження похідних даних;
\item зіставити заявлений тип, бінарну сигнатуру, граматичний розбір і результат декодування;
\item перевірити та використати сформовані $X_s$ і $X_m$ без виправлення виявлених відхилень;
\item обчислити структурні ознаки $Z_s$, оцінку $\rho_s$ і локалізації $L_s$;
\item сформувати з $X_v$ частотне представлення $X_f$, локальні та високочастотні характеристики;
\item для відео відібрати кадри, сформувати $X_t$, обчислити покадрові ознаки та виконати часове агрегування;
\item обчислити $Z_v$, оцінку $\rho_v$ і локалізації $L_v$;
\item сформувати статичний опис $H$ з окремими $Z_v$, $Z_s$, $L_v$, $L_s$, явною маскою й контекстом; похідне $Z_c=F_C(Z_v,Z_s)$ --- упорядковану конкатенацію двох канонічних векторів зі збереженням порядку компонентів і версії схеми ознак --- зберігати лише як числове поєднання.
\end{enumerate}

Перші два кроки є спільною умовою допустимості подальших обчислень. Після них структурна й спектрально-зорова гілки можуть виконуватися паралельно, але результат однієї не заповнює відсутнього результату іншої. Об’єднання на восьмому кроці дозволене лише після фіксації стану кожної гілки; завдяки цьому прискорення реалізації не змінює логіки методу.

Розгалуження та момент дозволеного об’єднання показано на рис.~\ref{fig:svs_method}. Стрілки від одного входу до двох гілок означають спільне походження, але не взаємозамінність їхніх результатів.

\begin{figure}[htbp]
\centering
\begin{tikzpicture}[x=1cm,y=1cm]
\node[dissertation block,text width=3.30cm] (input) at (0,3.15)
  {Статичні дані\\$X_c$};
\node[dissertation block,text width=3.70cm] (structural) at (-4.30,0.80)
  {Структурні ознаки\\$Z_s,L_s$};
\node[dissertation block,text width=3.70cm] (visual) at (0,0.80)
  {Детерміновані спектрально-зорові ознаки\\$Z_v,L_v$};
\node[dissertation block,dashed,text width=3.70cm] (learned) at (4.30,0.80)
  {Неперевірене розширення на основі навченої моделі};
\node[dissertation block,text width=8.10cm] (profile) at (0,-2.10)
  {Статичний опис\\$H=\langle Z_v,Z_s,\rho_v,\rho_s,L_v,L_s,m_c,C_H\rangle$};

\draw[dissertation arrow] (input.south) -- ++(0,-0.40) -| (structural.north);
\draw[dissertation arrow] (input.south) -- (visual.north);
\draw[dissertation arrow,dashed] (input.south) -- ++(0,-0.40) -| (learned.north);
\draw[dissertation arrow] (structural.south) -- ++(0,-0.45) -| (profile.north);
\draw[dissertation arrow] (visual.south) -- (profile.north);
\draw[dissertation arrow,dashed] (learned.south) -- ++(0,-0.45) -| (profile.north);
\end{tikzpicture}
\caption{Структура методу спектрально-зорового та структурного аналізу фото- і відеооб’єктів вебресурсів}
\label{fig:svs_method}
\end{figure}


Першим дослідницьким кроком є встановлення того, яке саме представлення дозволено аналізувати за правилами формату. Розширення, заявлений тип, сигнатура й поведінка засобу декодування можуть збігатися або суперечити одне одному, причому сама суперечність ще не доводить прихованого впливу. Окремо перевіряються чотири умови: відповідність розширення заявленому типу, заявленого типу --- бінарній сигнатурі, сигнатури --- установленому під час декодування типу, а також сам факт успішного декодування. Результати цих перевірок утворюють вектор $u_f$.

Роздільне збереження умов потрібне для пояснення причини подальшого маршруту. Наприклад, помилка декодування й наявність даних після логічного завершення контейнера потребують різних перевірок, хоча обидві можуть підвищити загальну структурну оцінку. Сам вектор $u_f$ не визначає прихованого впливу; він указує, де виникла неузгодженість і яку частину структури необхідно дослідити докладніше.

Після встановлення фактичного формату потрібно перейти від переліку контейнерних елементів до ознак, які можна порівнювати між об’єктами, не втрачаючи службових значень. Цей перехід конкретизує модельну функцію $F_s$ у формулі~\eqref{eq:ch3_structural_profile}:
\begin{equation}
Z_s=F_s(X_s,X_m),
\label{eq:ch3_structural_profile}
\end{equation}
де, $Z_s$ --- структурне представлення ознак; $X_s$ --- карта контейнера; $X_m$ --- представлення службових даних.

Функція $F_s$ не виправляє карти контейнера й не відкидає порушених службових значень. Вона переводить їх у зіставний опис, з якого надалі формується $\rho_s$, тоді як початкові байтові координати зберігаються в $L_s$. Отже, $Z_s$ впливає на глибину подальшої перевірки, а можливість пояснити цей вплив забезпечує саме зв’язок із $X_s$ та $L_s$.

Відтворюваність структурного аналізу залежить від того, чи однаково описано звичайний, вкладений і порушений елементи. Для кожного елемента $s_k$ зберігаються його тип $c_k$, початкова й кінцева байтові позиції $a_k$ та $b_k$, посилання $w_k$ на батьківський елемент і результат перевірки $v_k$. Такий склад потрібний не як ще одна математична модель, а як єдине правило реєстрації: без меж неможливо відтворити локалізацію, без батьківського елемента --- перевірити вкладеність, а без результату перевірки --- відрізнити прочитаний елемент від коректного.

Перевірка $v_k$ набуває додатного значення лише тоді, коли межі задовольняють умову $0\leq a_k<b_k\leq n$ і виконано форматні обмеження для відповідного типу, де $n$ --- довжина $X_{mm}$ у байтах. Поділ на дві умови необхідний через те, що фізично доступний інтервал може мати неправильний порядок або недопустиму вкладеність. Якщо заявлена довжина виходить за межі об’єкта, метод зберігає доступну частину, початкове значення поля та причину порушення. Байти після логічного завершення контейнера реєструються як окрема ділянка, а не відкидаються й не приписуються останньому елементу. Завдяки цьому подальший висновок спирається на фактично прочитані дані, а не на виправлену інтерпретацію контейнера.

\begin{table}[htbp]
\centering
\caption{Структурні перевірки мультимедійних форматів та очікувані результати розбору}
\label{tab:ch3_format_checks}
\small
\begin{tabularx}{\textwidth}{>{\raggedright\arraybackslash}p{0.12\textwidth}YYY}
\toprule
\textbf{Формат} & \textbf{Контрольовані структури} & \textbf{Ознаки неузгодженості} & \textbf{Результат} \\
\midrule
JPEG & Початок і кінець, службові сегменти, таблиці, скани & Порушені межі або порядок, дані після завершення & Карта сегментів і невіднесених ділянок \\
PNG & Сигнатура, основний заголовок, блоки даних, завершення, контрольні суми & Неможлива довжина, помилка контрольної суми, недопустиме повторення & Карта блоків і ознаки цілісності \\
WebP & Контейнер RIFF, блоки зображення, службові та анімаційні блоки & Розбіжність довжин, несумісний порядок, неврахований залишок & Карта контейнера й альтернативних інтерпретацій \\
Відео\-контей\-нери & Ієрархічні елементи, доріжки, часові шкали, індекси та потоки & Некоректна вкладеність, суперечлива тривалість, розрив часових міток & Ієрархічна карта й часові ділянки відхилень \\
\bottomrule
\end{tabularx}
\end{table}

Таблиця~\ref{tab:ch3_format_checks} не означає однакової повноти реалізації для всіх форматів. За наявними матеріалами контрольний цикл у розділі~\ref{chap:experiments} обмежено PNG і MP4. Для JPEG і WebP визначено правила розширення, але їх підтримку не слід вважати експериментально підтвердженою без відповідних аналізаторів і окремої перевірки.

Окремі формати містять різну кількість елементів і правил, тому абсолютна кількість спрацювань не є порівнюваною: одна помилка серед двох перевірок і одна серед ста мають різний зміст. Через це порушення спочатку об’єднуються у сім змістових груп: межі, порядок, службові дані, невіднесені ділянки, вкладені структури, багатоформатна інтерпретація та узгодженість декодування. Для кожної групи обчислюється частка порушених умов $d_g=n_g^{-}/\max(1,n_g)$, де $n_g^{-}$ --- кількість порушень, а $n_g$ --- кількість фактично виконаних перевірок. Нормування уможливлює зіставлення об’єктів різних форматів і не означає однакової важливості всіх правил.

Самої частки недостатньо, оскільки нуль може означати як відсутність порушень, так і неможливість виконати перевірку. Тому структурне представлення $Z_s$ містить вектор форматних розбіжностей $u_f$, сім часток $d_1,\ldots,d_7$ і вектор застосовності $q_s$. Якщо перевірку групи фактично виконано --- незалежно від того, виявила вона порушення чи ні, --- відповідна компонента $q_s$ дорівнює одиниці; якщо перевірка незастосовна до формату або недоступна --- нулю. Результат виконаних перевірок відображають частки $d_1,\ldots,d_7$, а не вектор застосовності. Завдяки цьому подальше обчислення не трактує невідоме значення як підтверджену норму, а виявлене порушення не змішується з невиконаною перевіркою. Порядок компонентів і версія схеми зберігаються разом з описом.

Повний $Z_s$ потрібний для пояснення висновку, тоді як вибір подальшого маршруту потребує однієї зіставної величини. Після формування вектора обчислюється структурна оцінка $\rho_s=\sigma(\beta_0+\beta^{\mathsf T}Z_s)$, де $\sigma(x)=1/(1+e^{-x})$ --- логістична функція, значення якої належать інтервалу $(0,1)$, а параметри $\beta_0$ і $\beta$ установлено на даних для налаштування. Ця оцінка не має імовірнісного змісту без окремого калібрування. Разом із нею зберігаються самі ознаки й відповідні байтові локалізації $L_s$; саме вони дають змогу з’ясувати, через яке порушення змінився маршрут аналізу.

Спектрально-зорова гілка розв’язує інше обмеження: слабке локальне відхилення може зникнути під час усереднення всього зображення або кадру. Через це $X_v$ поділяється на області однакового розміру та кроку, зафіксовані в $\Theta_H$. Для кожної області $r$ окремо визначаються частка одиниць у молодшій бітовій площині $b_r$, ентропія Шеннона локальної гістограми значень $e_r$, характеристика переходів між сусідніми значеннями $p_r$, міжканальна різниця $c_r$ і середній модуль високочастотного залишку $h_r$. Ці величини утворюють локальний опис $g_r$. Точні означення статистик $p_r$, $c_r$ і $h_r$ --- використані фільтри, канали та нормування --- зафіксовано у складі $\Theta_H$, а їх версія входить до контексту $C$.

Локальний опис не є самостійним рішенням щодо прихованого навантаження. Текстура, різкий край або сенсорний шум можуть створити подібний відгук без аномального коду. Призначення опису полягає у визначенні ділянок, для яких потрібно зберегти докладніше частотне представлення й контекст декодування. Остаточне тлумачення виконується лише після зіставлення з іншими областями, структурними ознаками та доступними подіями.

Локальні статистики не розкривають, у якому масштабі або частотному діапазоні виникло відхилення. Тому зорове представлення $X_v$ додатково перетворюється на частотні дані $X_f$. Для цього окремо застосовуються дискретне косинусне перетворення, дискретне хвилькове перетворення і залишкові фільтри; їхні рівні розкладу, частотні області та параметри утворюють набір $\theta_f\subset\Theta_H$. Роздільне збереження каналів необхідне через різне походження їхніх значень: підсумовування на цьому етапі не дозволило б установити, яке саме перетворення створило відгук.

Для JPEG початкові квантовані коефіцієнти, якщо вони доступні, зберігаються окремо від коефіцієнтів, повторно обчислених після декодування. Це обмеження усуває помилкове зіставлення величин, одержаних з різних станів об’єкта. Для кожної області та кожного частотного або залишкового каналу визначаються середній модуль коефіцієнтів, стандартне відхилення й ентропія гістограми. Їхній упорядкований набір утворює опис $u_r$.

Під час узагальнення потрібно одночасно врахувати поширеність слабкого відхилення й одиничний сильний прояв. Лише середнє послабило б коротку зміну, а лише найбільше значення зробило б опис надмірно чутливим до випадкового краю або шуму. Тому детерміноване представлення $z_d$ поєднує покомпонентні середні та найбільші значення локальних описів $g_r$ і частотних описів $u_r$. Наслідком такого рішення є збереження двох різних властивостей відхилення --- його поширеності та локальної виразності. Водночас $z_d$ залишається описом ознак, а не висновком про наявність аномального коду.

Можливе навчуване розширення повинно додатково зберігати положення ділянки, рівень частотного розкладу та формат носія, щоб ураховувати зв’язки між віддаленими областями. Для цього локальні описи можна впорядкувати за просторовими та масштабними індексами й подати навченому кодувальнику. За наявними матеріалами такий кодувальник не навчено й не перевірено незалежно. Отже, у прототипній конфігурації використовується лише $z_d$, а відсутній результат навчуваного розширення не замінюється нульовим вектором. Це обмеження не дає формальному опису майбутнього компонента набути статусу фактично одержаного результату.

Після локального та високочастотного аналізу потрібно повернутися до умов моделі й утворити один типізований вихід гілки, не приховуючи способу його одержання. Цей перехід конкретизує функцію $F_v$ у формулі~\eqref{eq:ch3_spectral_encoding}:
\begin{equation}
Z_v=F_v(X_v,X_f,X_t),
\label{eq:ch3_spectral_encoding}
\end{equation}
де, $Z_v$ --- спектрально-зорове представлення ознак; $X_t$ --- часове представлення кадрів; $F_v$ --- функція формування ознак.

В описаній прототипній конфігурації $F_v$ використовує детерміновані локальні характеристики, характеристики дискретного косинусного й дискретного хвилькового перетворень і залишкові характеристики. Додаткове навчуване представлення може ввійти до $Z_v$ лише після окремого навчання, фіксації параметрів та незалежної перевірки. Формула~\eqref{eq:ch3_spectral_encoding} задає спільний вихід методу й не приховує відмінності між фактично доступною та можливою розширеною конфігураціями.

Як і для структурної гілки, повний $Z_v$ зберігається для пояснення, а для вибору маршруту обчислюється проміжна оцінка. Для кожного опрацьованого кадру $j$ її задає логістичне перетворення $\rho_{v,j}=\sigma(\gamma_0+\gamma^{\mathsf T}Z_{v,j})$ з параметрами $\gamma_0$ і $\gamma$, налаштованими окремо від даних оцінювання; для зображення, яке містить єдиний кадр, $\rho_v=\rho_{v,1}$, а для відео агрегування покадрових оцінок визначає формула~\eqref{eq:ch3_frame_aggregation}. Покомпонентні значення $Z_v$ зберігаються разом із просторовими областями, оскільки сильний високочастотний відгук може бути наслідком текстури, межі, сенсорного шуму або повторного стискання. Через це $\rho_v$ впливає на глибину перевірки, але не використовується без форматного й локального контексту як самостійний висновок.

Відмінність між фактично описаним детермінованим формуванням ознак і необов’язковим навчуваним розширенням показано на рис.~\ref{fig:wavelet_tokenization}. Пунктирна гілка позначає лише напрям подальшого розвитку й не є свідченням виконаного експерименту.

\begin{figure}[htbp]
\centering
\begin{tikzpicture}[x=1cm,y=1cm]
\node[dissertation block,text width=3.20cm] (visual) at (-5.80,3.00)
  {Вхід за $m_v=1$\\
   $X_v,X_t,\theta_v$};
\node[dissertation block,text width=4.00cm] (transforms) at (-1.80,3.00)
  {Локальні характеристики, дискретне косинусне й хвилькове перетворення, залишкові фільтри\\
   $X_f$};
\node[dissertation block,text width=3.60cm] (local) at (2.45,3.00)
  {Локальні описи з просторово-часовими координатами};

\node[dissertation block,text width=3.60cm] (deterministic) at (-0.40,0.20)
  {Покомпонентні середні та найбільші значення\\
   $z_d$};
\node[dissertation block,text width=3.80cm] (output) at (-0.40,-2.40)
  {Підтверджена конфігурація\\
   $Z_v=z_d,\;L_v$};
\node[dissertation block,dashed,text width=3.70cm] (sequence) at (4.35,0.20)
  {Упорядкування локальних описів за простором, масштабом і часом};
\node[dissertation block,dashed,text width=3.90cm] (learned) at (4.35,-2.40)
  {Модель кодування не побудовано й не перевірено незалежно};

\draw[dissertation arrow] (visual.east) -- (transforms.west);
\draw[dissertation arrow] (transforms.east) -- (local.west);
\draw[dissertation arrow] (local.240) -- (deterministic.north);
\draw[dissertation arrow,dashed] (local.300) -- (sequence.north);
\draw[dissertation arrow] (deterministic.south) -- (output.north);
\draw[dissertation arrow,dashed] (sequence.south) -- (learned.north);
\end{tikzpicture}
\caption{Формування спектрально-зорових ознак у підтвердженій детермінованій конфігурації та межа неперевіреного розширення}
\label{fig:wavelet_tokenization}
\end{figure}


Для відео аналіз усіх кадрів збільшує витрати й породжує багато залежних спостережень, а аналіз лише ключових кадрів може пропустити короткий прояв. Тому відбір поєднує чотири підстави: належність до ключових кадрів, рівномірне покриття часу, межу сцени та близькість до структурно визначеної ділянки потоку. Об’єднання цих множин утворює $J$ без дублікатів. Крок рівномірного покриття, спосіб визначення меж сцен і радіус часової близькості до структурно визначеної ділянки належать до параметрів $\Theta_H$ і фіксуються у версії конфігурації. Такий спосіб зменшує обсяг обчислень, але зберігає кадри, важливі через часову або структурну причину.

Самого номера кадру недостатньо для відтворення перевірки. Тому в часовому представленні $X_t$ для кожного відібраного кадру $f_j$ зберігаються часова мітка $t_j$, усі причини відбору $u_j$ і показник технічної якості декодування $q_j$. Якщо кадр відповідає кільком умовам, він залишається одним записом, а причини накопичуються. Структурна причина включення не підвищує спектрально-зорової оцінки автоматично; вона лише гарантує, що відповідна часова ділянка не буде пропущена.

Покадрові спостереження мають неоднакову придатність: висока локальна оцінка в технічно пошкодженому кадрі не повинна мати той самий вплив, що й така сама оцінка за коректного декодування. Початкова значущість кадру обчислюється як зважена сума $a_j=\lambda_\rho\rho_{v,j}+\lambda_q q_j+\lambda_c u_{j,c}$, де $\rho_{v,j}$ --- покадрова оцінка; $q_j\in[0,1]$ --- показник якості декодування; $u_{j,c}\in\{0,1\}$ --- наявність структурної підстави відбору; $\lambda_\rho$, $\lambda_q$ і $\lambda_c$ --- невід’ємні параметри, встановлені на даних для налаштування. Невід’ємність усіх доданків гарантує $a_j\geq 0$, тому нормування значущостей до ваг $\omega_j=a_j/\sum_{l}a_l$ коректне; за нульової суми значущостей кадри одержують однакові ваги. Нормування потрібне, щоб кількість відібраних кадрів сама по собі не змінювала масштабу результату.

Після відбору й нормування відео має отримати один векторний опис та одну проміжну оцінку, причому обидва результати повинні спиратися на однакові ваги. Цю залежність визначено формулою~\eqref{eq:ch3_frame_aggregation}:
\begin{equation}
Z_v=\sum_{j=1}^{n_v}\omega_jZ_{v,j},\qquad
\rho_v=\sum_{j=1}^{n_v}\omega_j\rho_{v,j},
\label{eq:ch3_frame_aggregation}
\end{equation}
де, $n_v$ --- кількість різних відібраних кадрів; $Z_{v,j}$ --- представлення ознак кадру; $\rho_{v,j}$ --- його проміжна оцінка; $\omega_j$ --- нормована вага кадру; $Z_v$ і $\rho_v$ --- відповідно агреговане представлення та оцінка відео.

Спільні ваги у формулі~\eqref{eq:ch3_frame_aggregation} не дозволяють векторному результату й числовому маршруту спиратися на різні правила узагальнення. Оскільки покадрові оцінки $\rho_{v,j}$ належать інтервалу $(0,1)$, а ваги невід’ємні з одиничною сумою, агрегована оцінка $\rho_v$ також залишається в цьому інтервалі. Водночас велика вага кадру не є доказом аномального коду: вона лише визначає його внесок за зафіксованими умовами відбору та якості.

Для фото $n_v=1$, тому $\omega_1=1$. Параметри $\lambda_\rho$, $\lambda_q$ і $\lambda_c$ установлюються на даних для налаштування та не змінюються під час оцінювання.

Агрегована оцінка відповідає на питання про стан відео загалом, але не показує, у якій часовій ділянці виникла підстава рішення. Тому функція $G_v$ формує локалізації $L_v$ з $X_v$, $X_t$ і локального порога $\tau_v$. Спочатку залишаються кадри з $\rho_{v,j}\geq\tau_v$, далі вони впорядковуються за часом. Об’єднуються лише сусідні за часом кадри однієї сцени; граничний проміжок об’єднання зафіксовано серед параметрів $\Theta_H$. Для кожної групи зберігаються початкова й кінцева часові мітки, сукупність фактично спостережених областей і найбільша локальна оцінка. Неспостережений проміжок не оголошується підтвердженою локалізацією, оскільки для нього немає вихідного спостереження.

Етапи методу та їхні виходи узагальнено в таблиці~\ref{tab:ch3_svs_stages}.

\begin{table}[htbp]
\centering
\caption{Етапи методу спектрально-зорового та структурного аналізу, їхні операції та результати}
\label{tab:ch3_svs_stages}
\small
\begin{tabularx}{\textwidth}{>{\raggedright\arraybackslash}p{0.18\textwidth}YY}
\toprule
\textbf{Етап} & \textbf{Операції} & \textbf{Результат} \\
\midrule
Контроль входу & Перевірка ідентифікатора, контрольної суми й походження & Узгоджений $X_c$ \\
Форматна перевірка & Зіставлення заявленого типу, сигнатури, граматики й засобу декодування & $u_f$ \\
Структурна перевірка & Перевірка типів, меж, вкладеності та службових значень & $L_s$ \\
Структурне кодування & Нормування порушень і врахування застосовності & $Z_s$, $\rho_s$ \\
Локальний аналіз & Молодші біти, ентропія, переходи, міжканальні різниці, залишки & $g_r$ \\
Частотний аналіз & Дискретне косинусне перетворення, дискретне хвилькове перетворення, залишкові канали та прив’язка до областей & $X_f$, $z_d$ \\
Оброблення відео & Відбір, формування $X_t$, покадровий аналіз, зважування й локалізація & $X_t$, $Z_v$, $\rho_v$, $L_v$ \\
Формування опису & Інтеграція двох статичних гілок із маскою та контекстом & $H$ \\
\bottomrule
\end{tabularx}
\end{table}

Обчислювальну складність детермінованої конфігурації методу визначено відносно обсягу фактично опрацьованих представлень. Нехай $n_b$ --- кількість байтів початкового мультимедійного об’єкта, $n_p$ --- сумарна кількість піксельних позицій у відібраних кадрах, $n_e$ --- кількість структурних елементів контейнера, а $n_z$ --- кількість ознак, переданих до правила рішення. За послідовного розбору контейнера структурна гілка має порядок зростання часових витрат $O(n_b)$ і потребує до $O(n_e)$ пам’яті для структурної карти. Для спектрально-зорової гілки з фіксованими розмірами локальних блоків порядок зростання часових витрат становить $O(n_p)$; консервативна верхня оцінка пам’яті дорівнює $O(n_p)$, якщо декодовані кадри зберігаються одночасно. Формування спільного статичного опису має порядок $O(n_b+n_p)$, а застосування лінійного правила рішення --- $O(n_z)$. Ці оцінки характеризують залежність від розміру входу, але не враховують сталих витрат конкретного кодека, засобу декодування та програмної реалізації. Фактичні часові витрати вимірюються окремо в підрозділі~\ref{sec:ch4_ind_svs}.

Послідовність операцій показує, що результат методу визначається не лише значеннями $\rho_s$ і $\rho_v$. За доступності обох гілок утворюється повний статичний опис; за доступності однієї --- неповний. Відмову фіксують, якщо не підтверджено незмінності входу або походження даних. Відмова й неповнота не перетворюються на висновок про безпечність, оскільки в такому разі відсутня саме підстава для частини тверджень.

Отже, удосконалення методу полягає не в авторстві відомих дискретних косинусних, дискретних хвилькових або статистичних перетворень, а в узгодженій двогілковій процедурі, яка зберігає байтові й просторово-часові локалізації та формує один статичний опис без втрати семантики джерел. Одержаний $H$ є необхідним проміжним результатом: він дозволяє зіставити статичне спостереження з подією, але сам не доводить прояву прихованого впливу. Межі методу визначаються підтримкою форматної граматики, якістю декодування, параметрами кадрового відбору й відсутністю окремо перевіреного навчуваного зорового кодувальника. Незалежну перевірку детермінованої конфігурації для PNG і MP4 наведено в підрозділі~\ref{sec:ch4_ind_svs}. Саме потреба встановити зв’язок між статичним описом і реакцією середовища зумовлює другий метод.

\section{Метод поведінкового узгодження ознак різного походження}
\label{sec:behavior_alignment}

Статичний опис установлює, де й у якому представленні спостерігається відхилення, але не показує, чи пов’язане воно з реакцією компонента веб-системи. Журнал оброблення, навпаки, містить дії процесів, але без перевірки походження не дає змоги відрізнити наслідок поточного об’єкта від фонового навантаження. Безпосереднє приєднання журналу до $H$ створило б хибний причинний зв’язок, а відмова від подій залишила б невстановленим можливий прояв прихованого впливу.

Метою другого методу є інтеграція статичних ознак із подіями оброблення того самого мультимедійного об’єкта за підтвердженої належності й фазової сумісності. Дослідницьке припущення полягає в тому, що така перевірка дає змогу відокремити просту часову одночасність від змістовної підтримки статичної гіпотези. Метод спочатку встановлює причинну належність подій, потім формує $Z_b$ і лише після цього виконує інтеграцію. За відсутності придатного журналу статичний опис зберігається як неповний результат, а не доповнюється штучною нульовою поведінкою.

Оскільки інтегрувати потрібно три групи ознак, а не готові рішення окремих детекторів, вихід методу безпосередньо конкретизує центральну функцію моделі у формулі~\eqref{eq:ch3_alignment_mapping}:
\begin{equation}
Z=F_I(H,Z_b,m_A,C),
\label{eq:ch3_alignment_mapping}
\end{equation}
де, $F_I$ --- функція інтеграції ознак; $H$ --- статичний опис з окремими групами $Z_v$ і $Z_s$ та їх локалізаціями; $Z_b$ --- подієво-поведінкові ознаки; $m_A$ --- маска доступності; $C$ --- контекст походження та якості; $Z$ --- інтегроване представлення.

Канонічні $Z_v$ і $Z_s$ надходять до формули~\eqref{eq:ch3_alignment_mapping} як окремі складники опису $H$, а не відновлюються з похідного $Z_c$. Маска $m_A$ і контекст $C$ є явними аргументами, а локалізації $L_s$ і $L_v$ містяться в $H$; тому вихід функції узгоджується зі структурою $Z=\langle z,L,m_A,C\rangle$ і не втрачає підстав інтегрованого числового вектора.

На рис.~\ref{fig:behavior_alignment} суцільними лініями показано правилову (засновану на явних правилах) гілку, для якої в контрольному описі наявний агрегований результат, а пунктирними --- неперевірену необов’язкову гілку $U_b\rightarrow A$, доступну лише за виконання умов $m_A$ і $C$. Їхні результати надходять до різних входів $F_I$ без взаємного підмінювання. Таке розділення потрібне, щоб формалізація можливого розширення не сприймалася як доказ його фактичної реалізації.

\begin{figure}[htbp]
\centering
\begin{tikzpicture}[x=1cm,y=1cm]
\node[dissertation block,text width=3.25cm] (profile) at (-4.80,3.45)
  {Статичний опис\\$H$};
\node[dissertation block,text width=3.25cm] (events) at (0,3.45)
  {Початковий журнал\\$J_b$};
\node[dissertation block,text width=3.25cm] (context) at (4.80,3.45)
  {Доступність і супровідні\\відомості $m_A,C_I$};
\node[dissertation block,text width=3.25cm] (rules) at (-2.60,1.15)
  {Причинний відбір і явні\\правила $X_b,Z_r$};
\node[dissertation block,dashed,text width=3.25cm] (states) at (2.60,1.15)
  {Неперевірені послідовнісні стани\\$U_b$};
\node[dissertation block,dashed,text width=3.65cm] (attention) at (3.00,-1.05)
  {Неперевірене перехресне зіставлення $A$; до дослідного виходу не входить};
\node[dissertation block,text width=3.25cm] (integration) at (0,-2.85)
  {Інтеграція ознак\\$F_I$};
\node[dissertation block,text width=3.25cm] (joint) at (0,-4.65)
  {Інтегроване представлення\\$Z$};

\draw[dissertation arrow] ([xshift=-0.65cm]events.south)
  -- ++(0,-0.45) -| (rules.north);
\draw[dissertation arrow,dashed] ([xshift=0.65cm]events.south)
  -- ++(0,-0.45) -| (states.north);
\draw[dissertation arrow,dashed] (states.south) -- (attention.north);
\draw[dissertation arrow] (profile.south) -- ++(0,-0.45) -| (integration.west);
\draw[dissertation arrow,dashed] (context.east) -- ++(0.45,0) |- (attention.east);
\draw[dissertation arrow] (context.west) -- ++(-0.45,0) |- (integration.east);
\draw[dissertation arrow] (rules.south) -- ++(0,-0.45)
  -| ([xshift=-0.65cm]integration.north);
\draw[dissertation arrow] (integration.south) -- (joint.north);
\end{tikzpicture}
\caption{Підтверджена та неперевірена гілки формалізованої процедури поведінкового узгодження}
\label{fig:behavior_alignment}
\end{figure}


Відповідно до цієї межі спочатку встановлюється допустимість подій, потім їхнє представлення й лише наприкінці --- узгодження зі статичними даними:

\begin{enumerate}
\item відібрати події поточного запуску, що належать дозволеному ланцюгу процесів і стосуються об’єкта або простежуваної копії;
\item перевірити схему, обов’язкові поля, часовий порядок і покриття фаз;
\item канонізувати категорії та ресурси за версійованими словниками, зберігши посилання на початкові записи;
\item сформувати подієво-поведінкове представлення $X_b$;
\item обчислити правилове $Z_b$ або, за наявності навченої конфігурації, навчуване $Z_b$;
\item зіставити статичні локалізації з фазами та ресурсами причинно пов’язаних подій;
\item сформувати $Z$, оцінку узгодженості й показник упевненості без підміни відсутньої події нульовою ознакою.
\end{enumerate}

Перші три кроки не виділяють ознак, а встановлюють, чи має журнал право впливати на висновок. Четвертий і п’ятий формують його типізоване подання, шостий перевіряє змістовну відповідність статичним локалізаціям, а сьомий завершує інтеграцію. Якщо процес припиняється до завершення третього кроку, поведінкова гілка вважається недоступною; частково канонізовані записи не передаються далі як спостереження.

Подія набуває змісту не лише через тип, а й через місце в причинно-часовому ланцюгу. Після відсікання фонових записів придатні події впорядковуються в послідовність $X_b=\langle e_1,\ldots,e_T\rangle$, де $T$ --- кількість подій. Якщо причинно пов’язаних подій не одержано або мінімальне фазове покриття не виконано, установлюється $m_b=0$. Перелік обов’язкових фаз відповідає компонентам маршруту --- приймання, структурний розбір, декодування, перетворення та збереження --- і фіксується конфігурацією методу разом зі словниками. Стан $m_b=0$ означає недоступність спостереження, а не штатну поведінку; це розмежування не дозволяє технічній неповноті журналу зменшити оцінку ризику.

Однакова дія в різних фазах або над різними ресурсами може мати різний зміст. Для кожної події $e_t$ зберігаються дія $a_t$, ініціатор $i_t$, ресурс $r_t$, результат $o_t$ і час $\tau_t$. Фаза оброблення, ідентифікатор процесу та посилання на причинного попередника входять до контексту події. Події різних фаз, ініціаторів або результатів не об’єднуються. Суміжні однотипні записи можна згорнути лише зі збереженням їх кількості, часових меж і посилань на початкові записи. Завдяки цьому скорочення журналу не руйнує причинного ланцюга.

Типізований запис придатний для перевірки людиною, але обчислювальне зіставлення потребує числового представлення $u_t$. Його утворює функція $E_b$ за версійованими словниками категорій, часовими масштабами й масками полів. Наприклад, тип дії та клас ресурсу кодуються цілими індексами відповідних словників, результат --- бінарною ознакою успішності, а час --- нормованим зсувом відносно початку запуску. Відсутнє поле позначається окремою маскою, а не значенням, яке може збігатися з реальною категорією. Словники та порядок полів залишаються незмінними в межах однієї версії методу, оскільки їх зміна змінює зміст кожної компоненти.

Окреме кодування подій не відображає залежностей між віддаленими фазами. Тому повна конфігурація може додатково формувати впорядковані послідовнісні стани $U_b$ з числових описів $u_1,\ldots,u_T$. Проте за наявними матеріалами відповідний навчуваний кодувальник не навчено й не перевірено на незалежному наборі причинно пов’язаних послідовностей подій. Отже, послідовнісні стани визначають лише можливе розширення і не впливають на зафіксований прототипний результат.

Незалежно від обраної конфігурації вихід поведінкової гілки має відповідати одному модельному типу. Перехід від причинно пов’язаної послідовності до такого подання задано функцією $F_b$ у формулі~\eqref{eq:ch3_behavior_encoding}:
\begin{equation}
Z_b=F_b(X_b),
\label{eq:ch3_behavior_encoding}
\end{equation}
де, $F_b$ --- функція формування подієво-поведінкових ознак; $Z_b$ --- їх представлення.

Єдиний тип $Z_b$ звільняє функцію інтеграції від залежності від внутрішнього способу одержання поведінкових ознак. Водночас разом із результатом зберігається ознака використаної конфігурації, тому правиловий вихід не можна видати за результат неперевіреного навчуваного кодувальника. Значення $Z_b$ впливає на інтегроване представлення лише за підтвердженої доступності поведінкової гілки.

У повній конфігурації $F_b$ може використовувати $U_b$. У контрольному описі скорочену прототипну конфігурацію подано як правилову. Вона застосовує версійований каталог допустимих і нетипових відповідностей між фазою, дією, ресурсом, результатом і причинним попередником. Нормовані лічильники та індикатори правил утворюють представлення $Z_r=F_r(X_b,\theta_r)$, де $\theta_r$ позначає версію каталогу.

Для фактично описаної правилової конфігурації прийнято $Z_b=Z_r$. Таке рішення зумовлене відсутністю перевіреного навчуваного кодувальника й не означає, що правиловий опис є рівноцінним можливому розширенню. Для відтворення потрібні канонічний $X_b$, словники, каталог правил, порядок компонентів і пороги. Заміна каталогу створює іншу конфігурацію та має відображатися у версії результату.

Можливе навчуване розширення має розв’язати іншу задачу, ніж правилова конфігурація: зіставити не два вже узагальнені вектори, а конкретну статичну локалізацію з причинно сумісною подією. Для цього статичні елементи, сформовані з $Z_v$, $Z_s$, $L_v$ і $L_s$, та послідовнісні стани $U_b$ мають бути приведені до спільної розмірності. Однаковий розмір векторів, однак, не є достатньою підставою для їх поєднання, оскільки числова подібність може виникнути між подіями різних запусків або несумісних фаз.

Тому механізм розширення передбачає причинно-фазову маску $M$. Вона дозволяє зіставляти лише пари, що належать одному об’єкту, одному запуску, дозволеному ланцюгу процесів і сумісним фазам. Для дозволених пар обчислюються нормовані ваги відповідності, а їх результат утворює контекст $A$. Якщо статичний елемент не має жодної дозволеної події, відповідний рядок не перетворюється на звичайний нульовий вектор, а позначається недоступним. Ця умова не дозволяє відсутності спостереження виглядати як підтверджена відсутність поведінкового прояву.

Після зіставлення дозволені контексти можуть бути узагальнені й поєднані зі статичним представленням з урахуванням маски $m_A$. Внесок поведінкових даних має зменшуватися за неповної послідовності подій і бути відсутнім, коли причинна належність не підтверджена. Саме маска, а не числове значення заповнювача, визначає можливість такого поєднання. За наявними матеріалами послідовнісний кодувальник, причинно-фазове зважування та навчуване поєднання не налаштовано й не перевірено. Цей механізм описує умови можливого розширення, але не входить до результатів прототипної перевірки.

У правиловій конфігурації використовується простежуване поєднання похідного статичного представлення $Z_c=F_C(Z_v,Z_s)$ з правиловим вектором $Z_r$. Початкові $Z_v$ і $Z_s$ при цьому не вилучаються. Якщо $m_b=0$, правиловий внесок не обчислюється, а результат зберігає статичні ознаки й явну позначку неповноти. Якщо $m_b=1$, правила визначають, які поведінкові компоненти можуть підтримати або послабити конкретне статичне спостереження. Наслідком є відтворюване пояснення внеску подій, але не автоматичне доведення програмної семантики прихованого фрагмента.

Інтегроване представлення зберігає можливість подальшої класифікації, але для пояснення внеску поведінкової гілки потрібна окрема проміжна оцінка $\rho_b$. Її обчислюють із $Z_b$ лише за $m_b=1$; відсутній канал не створює значення $\rho_b=0$ як свідчення норми. Це обмеження впливає на рішення безпосередньо: за недоступної поведінкової гілки підсумковий стан має містити ознаку неповноти, а не занижену оцінку.

Висока виразність подій ще не означає прояву через конкретний компонент, оскільки та сама послідовність може бути штатною на іншій фазі або маршруті. Тому оцінка прояву $Q_A=F_Q(Z_b,R_{web},C)$ зіставляє поведінкові ознаки зі складом компонентів $R_{web}$ і причинним контекстом $C$. У правиловій конфігурації вона обчислюється як зважена частка спрацьованих правил нетипової відповідності серед фактично застосовних для маршруту правил каталогу; ваги правил $\alpha_r$ зафіксовано у версії каталогу, тому $Q_A\in[0,1]$, а незастосовні правила на оцінку не впливають. Одержане $Q_A$ зберігається в контексті виходу, входить до $Z$ і надходить до $F_D$. За наявними матеріалами окремої експериментальної перевірки цієї оцінки не проведено, тому вона визначає склад результату методу, але не є підтвердженою кількісною перевагою.

Після об’єднання важливо розрізнити силу взаємної підтримки ознак і якість підстав для такого зіставлення. Оцінка узгодженості визначається як $q_I=F_U(Z_v,Z_s,Z_b)$ і обчислюється через проміжні оцінки доступних гілок: $q_I=1-\overline{\Delta}$, де $\overline{\Delta}$ --- середнє попарних модулів різниць між доступними оцінками $\rho_v$, $\rho_s$ і $\rho_b$; за єдиної доступної гілки оцінка узгодженості не обчислюється й позначається недоступною. Високе $q_I$ означає сумісність груп ознак у визначеному причинному контексті, але не є підсумковим ризиком і не доводить програмної семантики прихованого фрагмента.

Окремо визначається упевненість інтеграції $c_I=F_N(m_A,q_p,q_t)$, де $q_p\in[0,1]$ --- частка покритих обов’язкових фаз, а $q_t\in[0,1]$ --- частка подій із підтвердженим причинним попередником. У прототипній конфігурації прийнято $c_I=\mu_A\min(q_p,q_t)$, де $\mu_A$ --- частка доступних гілок за маскою $m_A$; мінімум обрано, щоб жоден зі складників покриття не компенсував іншого. Це розділення необхідне через те, що сильна відповідність кількох наявних подій може спиратися на неповний журнал. Отже, $q_I$ характеризує сумісність свідчень, а $c_I$ обмежує силу висновку; низька упевненість не зменшує статичного ризику автоматично.

Обидві оцінки зберігаються в контексті $C$ інтегрованого представлення та надходять до $F_D$ як діагностичні входи для підсумкової упевненості $c_A$. Вони не замінюють ні $Q_A$, ні інтегрального ризику $\rho_A$.

Алгоритм правилової конфігурації формалізовано, однак її фактичне повторення потребує каталогу правил, канонічних подій, словників, порядку компонентів і порогів; повний комплект цих матеріалів у наявному корпусі не зафіксовано. Повного набору причинно пов’язаних послідовностей подій, навченого послідовнісного кодувальника, налаштованого перехресного зіставлення й незалежного порівняння немає. Тому одержаним результатом підрозділу є сукупність методичних умов причинного допуску, кодування та інтеграції, а не підтверджена перевага повної конфігурації. Її кількісні властивості не заявлено. Така межа безпосередньо визначає подальшу роль інформаційної технології: вона може організувати застосування доступної конфігурації та зберегти її статус, але не перетворює формалізований компонент на експериментально перевірений.

\section{Інформаційна технологія виявлення аномальних кодів у фото- та відеооб’єктах веб-ресурсів}
\label{sec:information_technology}

Окреме виконання двох методів ще не визначає, коли об’єкт можна передати веб-системі, що робити за неповного результату та як відтворити підстави рішення після зміни параметрів. Технологічна проблема, отже, полягає в керуванні життєвим циклом недовіреного носія, а не лише в послідовному запуску обчислень. До завершення аналізу потрібно зберігати ізоляцію, після аналізу --- відокремити стан від дії політики, а після дії --- зареєструвати достатній контекст для повторної перевірки.

Інформаційна технологія використовує одну модель і два методи, але не змінює їхнього наукового змісту. Модель формує аналітичний стан $S_A$, перший метод --- статичний опис, а другий --- інтегроване представлення. Технологія керує ізоляцією, порядком виконання, вибором дії, реконструкцією та реєстрацією. Саме це розмежування дає змогу змінити політику реагування, не приписуючи моделі нового висновку, і повторити перевірку, не замінюючи початковий об’єкт реконструйованим кандидатом.

Щоб кожна зміна стану мала визначений вхід, вихід і умову завершення, технологію організовано рівно в п’ять етапів; їх склад збігається з оглядом процесу в підрозділі~\ref{sec:process_organization} і тут закріплюється як незмінний склад технології:

\begin{enumerate}
\item \textit{Попереднє оброблення} --- приймання, ідентифікація, незмінне збереження, структурний розбір і контрольоване декодування.
\item \textit{Виділення ознак} --- виконання першого методу й контрольоване формування $X_b$ окремим каналом.
\item \textit{Інтеграція ознак} --- виконання другого методу та формування $Z$.
\item \textit{Прийняття рішення} --- формування $S_A$ і вибір технологічної дії за чинною політикою.
\item \textit{Реєстрація результатів} --- цілісне збереження стану, дії, версій і походження.
\end{enumerate}

П’ятиетапну послідовність показано на рис.~\ref{fig:information_technology}. Шлюз, ізольоване сховище, контрольоване середовище та сховище записів забезпечують виконання відповідних операцій, але не утворюють додаткових етапів і не змінюють логіки переходів.

\begin{figure}[htbp]
\centering
\begin{tikzpicture}[x=1cm,y=1cm]
\node[dissertation block,text width=3.65cm] (s1) at (-4.70,1.75)
  {1. Попереднє оброблення\\$S_0\rightarrow S_1$\\$X_{mm}\rightarrow X_c$};
\node[dissertation block,text width=3.65cm] (s2) at (0,1.75)
  {2. Формування ознак\\$S_1\rightarrow S_2$\\$H,J_b$};
\node[dissertation block,text width=3.65cm] (s3) at (4.70,1.75)
  {3. Відбір та інтеграція\\$S_2\rightarrow S_3$\\$X_b,Z$};
\node[dissertation block,text width=4.15cm] (s4) at (2.40,-1.40)
  {4. Прийняття рішення та визначення припису\\$S_3\rightarrow S_4$\\$S_4=\langle S_A,D_A\rangle$};
\node[dissertation block,text width=4.15cm] (s5) at (-2.40,-1.40)
  {5. Реєстрація результатів\\$S_4\rightarrow S_5$\\$S_5=F_R(S_4,C_R)$};

\draw[dissertation arrow] (s1.east) -- (s2.west);
\draw[dissertation arrow] (s2.east) -- (s3.west);
\draw[dissertation arrow] (s3.south) -- (s4.north);
\draw[dissertation arrow] (s4.west) -- (s5.east);
\end{tikzpicture}
\caption{П’ять етапів інформаційної технології виявлення ознак аномального вмісту в мультимедійних даних вебсередовища}
\label{fig:information_technology}
\end{figure}


Фіксований склад етапів не дозволяє часовим або якісним показникам окремої операції помилково стати характеристикою всього процесу. Інформаційна технологія $T_A$ охоплює п’ять упорядкованих етапів $T_1,\ldots,T_5$: попереднє оброблення, виділення ознак, інтеграцію ознак, прийняття рішення та реєстрацію результатів. Пропуск будь-якого з них змінює зміст результату. Зокрема, стан після прийняття рішення ще не є завершеним технологічним результатом, доки не збережено відомостей для його відтворення.

Шлюз приймання, ізольоване сховище, контрольоване середовище, засіб реконструкції та сховище записів є забезпечувальними компонентами. Кожна межа між ними має ідентифікований вхід, типізований вихід, контрольну суму, версію та стан завершення.

Аналітичний результат не визначає дії без урахування контексту застосування: однаковий ризик може бути прийнятним для внутрішнього архіву й неприйнятним для публічного каналу. Щоб така відмінність не змінювала модельного стану заднім числом, правило політики відокремлено від $F_D$. Залежність технологічної дії від аналітичного стану й чинної версії політики визначено формулою~\eqref{eq:ch3_policy_action}:
\begin{equation}
D_A=F_A(S_A,\theta_p),
\label{eq:ch3_policy_action}
\end{equation}
де, $D_A$ --- технологічна дія; $F_A$ --- функція застосування політики; $\theta_p$ --- версійовані параметри політики; $S_A$ --- аналітичний стан.

Допустимими діями є допуск, передавання на експертну перевірку, реконструкція та блокування. За високого ризику й низької упевненості політика не повинна автоматично обирати допуск. Однаковий $S_A$ може спричинити різні дії в різних середовищах, але дія не змінює модельного висновку. Формула~\eqref{eq:ch3_policy_action}, отже, дає змогу змінювати правила реагування без повторного тлумачення вже сформованого аналітичного стану.

Після застосування політики утворюється стан четвертого етапу, але повний цикл ще не завершено до успішної реєстрації. Склад цього проміжного стану, у якому аналітичний результат і дія залишаються окремими полями, визначено формулою~\eqref{eq:ch3_output_state}:
\begin{equation}
S_4=\langle S_A,D_A\rangle,
\label{eq:ch3_output_state}
\end{equation}
де, $S_4$ --- стан після етапу прийняття рішення; $S_A$ --- аналітичний результат; $D_A$ --- дія політики.

Подане у формулі~\eqref{eq:ch3_output_state} розділення дає змогу окремо порівнювати зміну ознак, зміну параметрів моделі та зміну політики.

Збереження лише $S_4$ не дає змоги пояснити зміну результату після оновлення засобу декодування, схеми ознак або політики. Тому п’ятий етап приєднує до цього стану контекст відтворення й утворює зареєстрований технологічний результат за формулою~\eqref{eq:ch3_audit_record}:
\begin{equation}
S_5=F_R(S_4,C_R),
\label{eq:ch3_audit_record}
\end{equation}
де, $F_R$ --- функція реєстрації; $S_5$ --- зареєстрований технологічний результат; $S_4$ --- стан після прийняття рішення; $C_R$ --- контекст відтворення.

Контекст $C_R$ містить ідентифікатор і контрольну суму $X_{mm}$, версії засобів, схем ознак і політики, ідентифікатор запуску, маску $m_A$, локалізації $L_s$ і $L_v$, перевірюване посилання на $X_b$, часові межі та причини неповноти. Посилання на журнал не замінює самого журналу під час повторного обчислення.

Реєстрація завершує процес. Якщо запис не створено цілісно, дія не вважається виконаною: об’єкт залишається в ізоляції, а система може призначити повторне опрацювання, експертний розгляд або блокування. Перерваний запуск не перезаписує завершеного запису; повторне виконання одержує новий ідентифікатор запуску.

Контрольована реконструкція потрібна тоді, коли політика допускає збереження дозволеного зорового вмісту, але не довіряє початковій структурі контейнера. Вона створює новий мультимедійний об’єкт без невідомих невіднесених ділянок і заборонених службових полів; це не очищення початкового носія, а поява іншого кандидата. Відповідне перетворення встановлено формулою~\eqref{eq:ch3_reconstruction_candidate}:
\begin{equation}
X_{mm,j}=E_R(X_{v,i},\theta_R),\quad j\neq i,
\label{eq:ch3_reconstruction_candidate}
\end{equation}
де, $E_R$ --- функція реконструкції; $X_{v,i}$ --- дозволене зорове представлення початкової перевірки; $\theta_R$ --- опис цільового формату; $X_{mm,j}$ --- новий мультимедійний об’єкт.

Реконструкція у формулі~\eqref{eq:ch3_reconstruction_candidate} не замінює $X_{mm,i}$ і не надає кандидату наперед установленого статусу безпечності. Зв’язок із джерелом зберігається в контексті, але новий об’єкт одержує власний ідентифікатор. Для відео перетворення охоплює лише зорове представлення: аудіодоріжки, субтитри та службові поля доріжок до нового кандидата автоматично не переносяться, а їх склад визначається описом цільового формату $\theta_R$ і чинною політикою.

Оскільки реконструкція змінює байтову послідовність, попередній клас або ризик не можна успадковувати як властивість нового об’єкта. Кандидат повертається на початок і проходить повний технологічний цикл за формулою~\eqref{eq:ch3_reconstruction_cycle}:
\begin{equation}
S_{5,j}=F_T(X_{mm,j},\theta_T),
\label{eq:ch3_reconstruction_cycle}
\end{equation}
де, $F_T$ --- функція виконання п’яти етапів технології; $\theta_T$ --- зафіксована конфігурація; $S_{5,j}$ --- зареєстрований результат нового п’ятиетапного циклу.

Новий цикл у формулі~\eqref{eq:ch3_reconstruction_cycle} не успадковує класу, ризику або дії попередньої перевірки. Допуск можливий лише за результатом нового аналітичного стану й успішної реєстрації. Реконструкція при цьому не утворює шостого етапу технології: вона не продовжує поточної перевірки, а породжує новий об’єкт і, відповідно, новий п’ятиетапний цикл із власним ідентифікатором.

Ізоляцію об’єкта до зареєстрованого рішення та повернення лише реконструйованого кандидата на початок нового циклу показано на рис.~\ref{fig:zero_trust}. Зворотний зв’язок на схемі тому не означає повторного використання попереднього рішення.

\begin{figure}[htbp]
\centering
\begin{tikzpicture}[x=1cm,y=1cm]
\node[dissertation block,text width=4.20cm] (object) at (0,4.85)
  {Початковий об’єкт\\$X_{mm}$};
\node[dissertation block,text width=4.20cm] (isolation) at (0,3.25)
  {Ізоляція об’єкта};
\node[dissertation block,text width=4.20cm] (cycle) at (0,1.65)
  {Формування аналітичного стану\\$S_A$};
\node[dissertation block,text width=4.20cm] (registered) at (0,0.05)
  {Політика реагування\\$D_A=F_A(S_A,\theta_p)$};
\node[dissertation block,text width=4.20cm] (policy) at (0,-1.55)
  {Зареєстрований результат\\$S_5=F_R(S_4,C_R)$};

\node[dissertation block,text width=2.45cm] (allow) at (-5.10,-4.00)
  {Допуск};
\node[dissertation block,text width=2.85cm] (review) at (-1.75,-4.00)
  {Експертна перевірка};
\node[dissertation block,text width=2.45cm] (block) at (1.75,-4.00)
  {Блокування};
\node[dissertation block,text width=3.35cm] (reconstruct) at (5.15,-4.00)
  {Реконструкція кандидата\\$X_{mm,j}$};

\draw[dissertation arrow] (object.south) -- (isolation.north);
\draw[dissertation arrow] (isolation.south) -- (cycle.north);
\draw[dissertation arrow] (cycle.south) -- (registered.north);
\draw[dissertation arrow] (registered.south) -- (policy.north);

\coordinate (decisionbus) at (0,-2.80);
\draw[dissertation arrow] (policy.south) -- (decisionbus);
\draw[thick] (-5.10,-2.80) -- (5.15,-2.80);
\draw[dissertation arrow] (-5.10,-2.80) -- (allow.north);
\draw[dissertation arrow] (-1.75,-2.80) -- (review.north);
\draw[dissertation arrow] (1.75,-2.80) -- (block.north);
\draw[dissertation arrow] (5.15,-2.80) -- (reconstruct.north);

\coordinate (loopright) at (7.15,3.25);
\draw[thick] (reconstruct.east) -- (7.15,-4.00) -- (loopright);
\node[font=\small,anchor=east] at (6.95,-0.85) {новий цикл};
\draw[dissertation arrow] (loopright) -- (isolation.east);
\end{tikzpicture}
\caption{Контур нульової довіри для інформаційної технології}
\label{fig:zero_trust}
\end{figure}


Синхронний, асинхронний і фоновий режими не змінюють п’яти етапів. Перевищення часу чи пам’яті, помилка декодування та відмова сховища є технічними станами, а не нульовим ризиком. Об’єкт не публікується автоматично через технічну невдачу.

Алгоритмічна відтворюваність технології потребує незмінних входів, версійованих параметрів, канонічних схем подій, каталогу правил, порогів, ваг і повних контрольних записів. Однак формальна наявність цих вимог не доводить їх реалізації в розподіленому середовищі. За наявними матеріалами не перевірено роботу черги, цілісність переходів, відновлення після відмови, навантажувальні характеристики та повний час усіх п’яти етапів. Отже, сформовано послідовність і умови виконання інформаційної технології, але її промислова готовність не заявляється; експериментальний розділ може оцінювати лише той операційний зріз, для якого збережено фактичні дані.

\section*{Висновки до розділу 3}
\addcontentsline{toc}{section}{Висновки до розділу 3}

У розділі розв’язано задачу переходу від модельного опису до послідовності операцій зі збереженням походження кожного свідчення. Для цього відокремлено незмінний вхід від похідних представлень, статичний опис --- від причинно пов’язаного журналу, аналітичний стан --- від дії політики. Модель розділу~\ref{chap:models_formalization} визначає дані й аналітичний стан, перший метод формує статичний опис, другий --- інтегроване представлення, а інформаційна технологія організовує їх п’ятиетапне застосування та реєстрацію.

Удосконалено метод спектрально-зорового та структурного аналізу через двогілкове формування ознак зі збереженням байтових і просторово-часових локалізацій. Обраний порядок усуває втрату контейнерних свідчень під час декодування та передчасне змішування фізично різних ознак. Визначено вхід, параметри, вісім операційних кроків, правила завершення, структурні перевірки, детерміновані локальні й частотні ознаки, кадровий відбір та агрегування. Метод формує $H$ і не підміняє інтеграції або технологічного рішення. Детерміновану конфігурацію окремо перевірено за групово незалежним протоколом у підрозділі~\ref{sec:ch4_ind_svs}; кількісна перевага повної навчуваної конфігурації цим результатом не встановлюється.

Запропоновано метод поведінкового узгодження ознак різного походження, у якому проста часова одночасність не є достатньою підставою для інтеграції: подія повинна належати тому самому запуску, дозволеному причинному ланцюгу та сумісній фазі. Для описаної в контрольному описі скороченої конфігурації формалізовано правилове представлення з версійованим каталогом відповідностей. Навчувану конфігурацію з послідовнісним кодувальником і перехресним зіставленням відокремлено як неперевірене розширення. Повний комплект правил і послідовностей подій у наявному корпусі не зафіксовано, тому експериментальну перевагу методу та кількісні властивості повної конфігурації не заявлено.

Сформовано інформаційну технологію як послідовність попереднього оброблення, виділення ознак, інтеграції ознак, прийняття рішення та реєстрації результатів. Аналітичний стан $S_A$ відокремлено від технологічної дії $D_A$ у стані $S_4=\langle S_A,D_A\rangle$, а п’ятий етап утворює зареєстрований результат $S_5$. Реконструйований кандидат одержує новий ідентифікатор і проходить повний новий цикл.

Межі одержаних результатів визначаються двома різними циклами. Ретроспективне кількісне оцінювання в розділі~\ref{chap:experiments} стосується агрегованих результатів зафіксованої в контрольному описі конфігурації, повторення якої обмежене збереженими похідними показниками, тоді як незалежний експеримент у підрозділі~\ref{sec:ch4_ind_svs} перевіряє детерміновані структурні та спектрально-зорові ознаки на синтетичних PNG і MP4. Жоден із цих циклів не підтверджує повних навчуваних компонентів, причинно пов’язаних послідовностей подій, додаткових форматів або промислового контуру.
